\title{Fingerübung Heizen + Kühlen}
\author{Architektur und Technologie}
\date{}
\SemioAssignmentSetup{course={Architektur und Technologie},number={Fingerübung},cover-content={\EnergyCoverAxonometry},cover-title={Tiny Home in Rhino und Grasshopper},submission={CSV-Datei auf Stud.IP bis Freitag, 10:00 Uhr}}
\SemioHeaderSetup{left={Architektur und Technologie},center={Fingerübung},right={Heizen + Kühlen}}
\newcommand{\EnergyDecimal}{{,}}
\newcommand{\EnergyWordElement}{Element}
\newcommand{\EnergyWordEnvelope}{Gebäudehülle}
\newcommand{\EnergyWordFloor}{Nutzfläche}
\newcommand{\EnergyWordWindowSouth}{Fenster Süd}
\newcommand{\EnergyWordDoorEast}{Schiebetür Ost}
\newcommand{\EnergyWordAxonometry}{Axonometrie}
\newcommand{\EnergyWordCoordinates}{Koordinaten in Metern}
\newcommand{\EnergyWordCsvExample}{Aufbau der Zeilen}
\newcommand{\EnergyWordCsvRows}{Zeilen der Vorlage}
\newcommand{\EnergyWordCanopy}{Vordach}
\begin{document}
\makeassignmentcover
\section{Aufgabe}
Bestimme den Jahresheizwärmebedarf und die zugehörigen Heizenergiekosten dieses Tiny Homes mit einem kleinen Rhino-Grasshopper-Skript. Implementiere zusätzlich die Kühllast für die gläserne Schiebetür und das große Fenster im Süden und bestimme den erforderlichen Rohrdurchmesser einer Lüftungsanlage. Wir empfehlen dir, die gesamte Übung direkt in Grasshopper umzusetzen: Geometrie, Berechnung und Export in einer Definition.

\subsection{Vorgegeben}
Exakt einzuhalten sind nur zwei Dinge: die \textbf{Geometrie} (Abschnitt 2) und der \textbf{Aufbau der CSV-Datei} (Abschnitt 4). Zusätzlich gilt als Rahmen: Das Haus steht in Hannover, es wird mit einer Luftwärmepumpe beheizt und hat eine normale Isolierverglasung, also keine Sonnenschutzverglasung.

\subsection{Nicht vorgegeben}
Konkrete Energiewerte nennt dieses Blatt bewusst nicht; sie zu finden, ist Teil der Übung. Das betrifft alle Kennwerte der Bilanz: Klimadaten des Standorts, U-Werte der Bauteile, Luftwechsel und Luftkennwerte, das Verhältnis von Innen- zu Außenvolumen, solare Einstrahlung und Abminderungsfaktoren, interne Gewinne, Jahresarbeitszahl, Strompreis, Gesamtenergiedurchlassgrad, Auslegungseinstrahlung, Temperaturdifferenz der Zuluft und zulässige Luftgeschwindigkeit. Entnimm diese Werte der Vorlesung, den dort genannten Normen oder einer eigenen Recherche. Lege jeden Wert als benannte Eingabe an und halte die Quelle in der Definition fest.
\section{Geometrie}
Alle Angaben sind Koordinaten in Metern im Weltkoordinatensystem von Rhino: $x$ zeigt nach Osten, $y$ nach Norden, $z$ nach oben. Modelliere exakt diese Geometrie.
\EnergyAxonometryFigure
\EnergyElements
Die Gebäudehülle ist ein geschlossener Quader mit Flachdach. Das Fenster im Süden und die Schiebetür im Osten sind ebene Rechtecke in der Fassadenebene; ihre Maße sind Öffnungsmaße einschließlich Rahmen. West- und Nordfassade, Dach und Boden haben keine Öffnungen. Die Nutzfläche entspricht vereinfacht der Grundfläche.

Das Vordach ist eine ebene, waagerechte, opake Fläche über der Schiebetür. Es verschattet und gehört nicht zur thermischen Hülle.
\section{Berechnung}
Folge dem Rechenweg der Vorlesung „Heizen + Kühlen“. Leite alle Flächen und Volumen aus der Geometrie ab, statt sie abzutippen, und rechne intern ungerundet.
\subsection{Heizen}
\begin{enumerate}
\item Ermittle je Himmelsrichtung die opake Wandfläche und die verglaste Fläche sowie Dachfläche, Bodenfläche, Außenvolumen und Nutzfläche.
\item Berechne die Transmission jedes Bauteils und die Transmissionsverluste $H_T$, die Lüftungsverluste $H_V$ und daraus mit den Klimadaten die Wärmeverluste $Q_{l,HP}$.
\item Berechne die solaren Gewinne $Q_{s,HP}$ je Himmelsrichtung und die internen Gewinne $Q_{i,HP}$.
\item Bilanziere den Jahresheizwärmebedarf $Q_h$ und den flächenbezogenen Wert $q_h$.
\item Bestimme für die Luftwärmepumpe den Endenergiebedarf $Q_E$ und die jährlichen Energiekosten.
\end{enumerate}
\subsection{Kühlen}
\begin{enumerate}
\item Lege für die Schiebetür und das Fenster im Süden den sommerlichen Auslegungsfall fest und begründe ihn. Entscheide, wie du die besonnte Fläche $A$ ansetzt und ob das Vordach sie verringert.
\item Berechne die Kühllast $\dot Q_{S,tr}$ aus Fläche, Gesamtenergiedurchlassgrad $g_{tot}$ und Einstrahlung $I_S$.
\item Bestimme aus der Temperaturdifferenz $\Delta\Theta$ den erforderlichen Volumenstrom $q_{v,R}$, daraus den Rohrquerschnitt $A_{Rohr}$ und den Durchmesser $d$ eines runden Rohres.
\end{enumerate}
\subsection{Empfehlungen für Grasshopper}
\begin{enumerate}
\item Baue die Geometrie parametrisch auf oder referenziere sie aus Rhino, und miss Flächen und Volumen mit Komponenten wie \emph{Area}, \emph{Volume} und \emph{Deconstruct Brep}.
\item Benenne Eingaben und Ergebnisse mit Größe und Einheit. Trage keine Ergebnisse von Hand ein; jede Ausgabe muss auf Änderungen der Eingaben reagieren.
\item Gruppiere die Definition nach Geometrie, Heizen, Kühlen und Export. Für den Export genügen native Komponenten oder ein kurzes Skript ohne zusätzliches Plug-in.
\end{enumerate}
\section{Abgabe}
Lade genau eine Datei \EnergySubmissionFile{} bis Freitag, 10:00 Uhr, auf Stud.IP in den Ordner \emph{\EnergyStudIpUpload} hoch; ersetze NAME durch deinen Namen.
\begin{itemize}
\item Vorlage: \EnergyStudIpFiles{} / \EnergyTemplateFile
\item Referenzlösung: \EnergyStudIpFiles{} / Referenz / \texttt{energiebilanz.csv}, erstellt mit der Definition \texttt{entwurf.gh} im selben Ordner. Die Referenz zeigt Aufbau und Vorgehen an einem anderen Entwurf; ihre Zahlen gelten nicht für dieses Tiny Home.
\end{itemize}
Der Aufbau deiner Datei muss exakt der Vorlage entsprechen: 60 Zeilen mit je zwei Feldern, Bezeichnung und Wert.
\EnergyCsvExample{Wert}
\begin{enumerate}
\item Kodierung UTF-8 mit BOM, Zeilenende CRLF, Trennzeichen Semikolon, jedes Feld in doppelten Anführungszeichen.
\item Bezeichnungen, Reihenfolge und Leerzeilen bleiben unverändert. Kopiere die Bezeichnungen aus der Vorlage, statt sie abzutippen; einzelne enden mit einem Leerzeichen, Zeile 41 beginnt mit dem Auslassungszeichen.
\item Überschriften (zum Beispiel Zeile 1, 5 und 18) und Leerzeilen erhalten einen leeren Wert.
\item Zahlen stehen mit Dezimalkomma, ohne Tausendertrennzeichen und ohne Einheit; die Einheit ergibt sich aus der Bezeichnung. Runde erst beim Export, auf höchstens vier Nachkommastellen.
\item Zeile 44 enthält die Art der Wärmepumpe als Text. Bauteile, die es nicht gibt, erhalten den Wert 0.
\end{enumerate}
\EnergyCsvTemplate

\small\textit{Inhaltliche Grundlage: Vorlesung „Heizen + Kühlen“, Nachhaltige Gebäudesysteme, Leibniz Universität Hannover.}
\end{document}
